% =====================================================================
% Artigo técnico — Projeto Final GraphRAG & KAG
% Template: SBC Conferences (sbc-template.sty + sbc.bst, disponíveis no
% template "SBC Conferences" do Overleaf).
%
% Cada seção fica em um arquivo próprio em secoes/, com o responsável
% indicado no topo. Cada pessoa edita SÓ o seu arquivo, para evitar
% conflitos no Git. Trechos ainda não escritos aparecem em vermelho no PDF.
% =====================================================================

\documentclass[12pt]{article}

\usepackage{sbc-template}
\usepackage{graphicx,url}
\usepackage[utf8]{inputenc}
\usepackage[T1]{fontenc}
\usepackage[brazil]{babel}
\usepackage{booktabs}
\usepackage{tabularx}
\usepackage{xcolor}
\usepackage{tikz}
\usetikzlibrary{positioning,arrows.meta,fit,backgrounds}

% Marca trechos que ainda precisam ser escritos (aparece em vermelho).
% Antes da entrega final, nenhum \pendente pode sobrar.
\newcommand{\pendente}[1]{\textcolor{red}{\textbf{[PENDENTE: #1]}}}

% Nomes de rótulos e relações do grafo em fonte monoespaçada.
\newcommand{\rot}[1]{\texttt{:#1}}
\newcommand{\rel}[1]{\texttt{#1}}

\sloppy

\title{Recuperação Aumentada por Grafos de Conhecimento\\
sobre a História da Computação}

\author{Maria Eduarda Zumbé\inst{1}, Anna Maria Albino\inst{1},
  Pessoa 3\inst{1}, Pessoa 4\inst{1}, Pessoa 5\inst{1}, Pessoa 6\inst{1}}

\address{Curso de Bacharelado em Ciência da Computação\\
  \pendente{nome da instituição -- cidade -- UF -- Brasil}
  \email{\pendente{e-mails dos autores}}
}

\begin{document}

\maketitle

% =====================================================================
% Abstract e Resumo
% Responsável: Pessoa 1 (Maria Eduarda) — escrever POR ÚLTIMO, depois
% que os resultados estiverem prontos.
%
% Cada um com até 10 linhas: problema, o que foi feito, principal
% resultado (com número) e conclusão.
% =====================================================================

\begin{abstract}
\pendente{abstract in English, written last}
\end{abstract}

\begin{resumo}
\pendente{resumo em português, escrito por último}
\end{resumo}


% =====================================================================
% 1. Introdução
% Responsável: Pessoa 1 (Maria Eduarda)
%
% O que cobrir (cerca de 1 página):
% - Problema: LLMs alucinam, falham em raciocínio multi-hop e não
%   mostram de onde tiraram a resposta.
% - Por que grafos ajudam: conhecimento explícito, travessia de
%   relações, procedência verificável.
% - O que este trabalho faz: sistema GraphRAG/KAG sobre a história da
%   computação, com Neo4j + LLM.
% - Contribuições, em lista: (i) dataset, (ii) modelo LPG, (iii)
%   pipeline Text-to-Graph, (iv) recuperação híbrida multi-hop,
%   (v) interface com procedência, (vi) avaliação comparativa.
% - Parágrafo final com a organização do artigo (seção 2 trata de...).
% =====================================================================

\section{Introdução}

\pendente{introdução -- escrever depois que as seções técnicas estiverem prontas}

% =====================================================================
% 2. Fundamentação Teórica e Trabalhos Relacionados
% Responsável: Pessoa 4
%
% O que cobrir (cerca de 1 página):
% - RAG tradicional (busca vetorial + LLM) e suas limitações
%   \cite{lewis2020rag}.
% - Grafos de conhecimento e o modelo de grafo de propriedades
%   \cite{hogan2021kg,angles2018pg}.
% - GraphRAG \cite{edge2024graphrag} e KAG \cite{liang2024kag}: o que
%   cada um propõe e em que este trabalho se apoia.
% - Diferença deste trabalho: domínio, idioma (português) e foco em
%   procedência visível na interface.
% As referências acima já estão em referencias.bib.
% =====================================================================

\section{Fundamentação Teórica e Trabalhos Relacionados}

\pendente{fundamentação teórica e trabalhos relacionados}

% =====================================================================
% 3. Dataset
% Responsável: Pessoa 1 (Maria Eduarda) — ESCRITA
%
% Números calculados a partir de data/raw/_relatorio.csv. Se a coleta
% for refeita, atualizar a Tabela 1 e os totais do texto.
% =====================================================================

\section{Dataset}
\label{sec:dataset}

O domínio escolhido foi a história da computação: as pessoas, organizações, linguagens de programação, máquinas e tecnologias que marcaram a área. O domínio foi escolhido por três motivos. Primeiro, os textos sobre o tema são narrativos e não estruturados, o que exige uma etapa real de extração de entidades e relações. Segundo, as mesmas entidades aparecem em muitos documentos (o Bell Labs, por exemplo, é citado nas páginas de pessoas, linguagens e sistemas operacionais), o que gera conexões entre documentos e torna necessária a resolução de entidades. Terceiro, o domínio favorece perguntas que só podem ser respondidas combinando várias relações, como ``que linguagens foram criadas por pessoas que trabalharam no Bell Labs?'', que é o tipo de consulta multi-hop que motiva o uso de grafos.

\subsection{Seleção das páginas}

O corpus é formado por artigos da Wikipédia em português. A escolha por um único idioma, o mesmo das perguntas dos usuários, evita que a busca textual e a resolução de entidades precisem lidar com nomes e termos em línguas diferentes.

As páginas não foram escolhidas ao acaso. A seleção partiu de 20 páginas \emph{semente}, centrais na história da computação (por exemplo, Alan Turing, Bell Labs, Unix, linguagem C, ENIAC, ARPANET e transistor), e foi ampliada com páginas citadas por elas. Com isso, o corpus se organiza em grupos densamente conectados, como o do Bell Labs (Dennis Ritchie, Ken Thompson, Unix, C, PDP-11), o que vai do transistor ao microprocessador (Shockley, Fairchild, Intel, Intel 4004, Altair 8800) e o da Internet (Vint Cerf, Robert Kahn, TCP/IP, ARPANET, CERN, World Wide Web). As páginas foram distribuídas em cinco categorias, apresentadas na Tabela~\ref{tab:dataset}.

\begin{table}[ht]
\centering
\caption{Composição do dataset por categoria.}
\label{tab:dataset}
\begin{tabular}{lrr}
\toprule
\textbf{Categoria} & \textbf{Páginas} & \textbf{Caracteres (mil)} \\
\midrule
Pessoa                     & 40  & 483   \\
Organização                & 29  & 647   \\
Linguagem de programação   & 29  & 433   \\
Hardware                   & 30  & 199   \\
Sistema, invenção ou evento & 19 & 421   \\
\midrule
\textbf{Total}             & \textbf{147} & \textbf{$\approx$\,2.180} \\
\bottomrule
\end{tabular}
\end{table}

A lista inicial tinha 151 páginas. Quatro foram removidas por não existirem na Wikipédia em português (Centrum Wiskunde \& Informatica, a linguagem Logo, o navegador Mosaic e a Conferência de Dartmouth). % CONFERIR: buscar "Dartmouth" nos arquivos de McCarthy e Minsky em
% data/raw/. Se não aparecer, apagar o trecho entre parênteses.
Essas remoções tiveram pouco impacto nas conexões do corpus, pois os temas correspondentes também são tratados em outras páginas (por exemplo, a Conferência de Dartmouth está ligada a John McCarthy e Marvin Minsky, cujas páginas fazem parte do dataset). Uma página que apontava para uma página de desambiguação (Colossus) foi substituída pela página do computador correspondente.

\subsection{Coleta e pré-processamento}

Os textos foram coletados em outubro de 2026 pela API do MediaWiki, por um script em Python disponível no repositório do projeto. Para cada página, a API retorna o texto em formato simples, o título final (após seguir redirecionamentos), a URL e o identificador da página no Wikidata. Doze títulos da lista foram resolvidos por redirecionamento, como ``Transistor'' para ``Transístor'' e ``Mouse'' para ``Rato (informática)''.

No pré-processamento, foram removidas as seções finais dos artigos que não descrevem a entidade (Referências, Ligações externas, Bibliografia, Ver também e Notas), e as quebras de linha repetidas foram normalizadas. Cada página foi salva em um arquivo JSON com o texto limpo e seus metadados: título, URL, identificador do Wikidata, categoria, data da coleta e licença.

O corpus resultante tem cerca de 2,2 milhões de caracteres, com mediana de aproximadamente 10 mil caracteres por página. O tamanho das páginas varia muito: de 387 caracteres (Robert Kahn) a mais de 110 mil (John von Neumann). Dezessete páginas têm menos de 1.500 caracteres. % CONFERIR: buscar "Kahn" em data/raw/TCP_IP.json e data/raw/ARPANET.json
% e ajustar o exemplo conforme o que for encontrado.
Elas foram mantidas de propósito, pois essas entidades também são citadas em páginas maiores do corpus: Robert Kahn, por exemplo, está associado ao desenvolvimento do TCP/IP e da ARPANET, ambos presentes no dataset. Combinar informações espalhadas por várias páginas é justamente o que a travessia multi-hop no grafo deve explorar.

\subsection{Licença e uso para avaliação}

O conteúdo da Wikipédia é distribuído sob a licença CC~BY-SA~4.0, o que permite seu uso neste trabalho com a devida atribuição. Além do texto, cada página traz o identificador do Wikidata, base de conhecimento estruturada mantida pela mesma comunidade. Esses identificadores permitem comparar as relações extraídas automaticamente com fatos curados do Wikidata, o que é usado na avaliação descrita na Seção~\ref{sec:avaliacao}.

% =====================================================================
% 4. Modelagem do Grafo
% Responsável: Pessoa 1 (Maria Eduarda) — ESCRITA
%
% Baseada em docs/modelo_lpg.md e schema/schema.cypher. Se o modelo
% mudar, atualizar as tabelas e a figura.
% =====================================================================

\section{Modelagem do Grafo}
\label{sec:modelagem}

O conhecimento extraído é armazenado no Neo4j segundo o modelo de grafo de propriedades rotuladas (\emph{Labeled Property Graph}, LPG) \cite{angles2018pg,robinson2015graphdb}. Nesse modelo, nós e relações têm um tipo (rótulo) e podem carregar propriedades na forma chave--valor, o que permite representar tanto as entidades do domínio quanto informações sobre a origem de cada fato.

\subsection{Duas camadas: conhecimento e procedência}

O grafo foi dividido em duas camadas, ilustradas na Figura~\ref{fig:modelo}. A \emph{camada de conhecimento} contém as entidades do domínio e as relações entre elas; é sobre ela que ocorre a travessia multi-hop. A \emph{camada de procedência} contém os documentos da Wikipédia (\rot{Documento}) e os trechos em que foram divididos (\rot{Chunk}). As duas camadas se ligam pela relação \rel{MENCIONA}, de um trecho para cada entidade citada nele, e pela relação \rel{SOBRE}, de um documento para a entidade de que ele trata.

Essa separação atende diretamente ao objetivo de reduzir alucinações: qualquer entidade ou relação usada para responder a uma pergunta pode ser rastreada até o trecho original de onde foi extraída, e esse trecho pode ser mostrado ao usuário como evidência.

\begin{figure}[ht]
\centering
\resizebox{\linewidth}{!}{%
\begin{tikzpicture}[
  ent/.style={draw, rounded corners=3pt, minimum width=2.2cm, minimum height=0.7cm, font=\footnotesize},
  prov/.style={ent, fill=gray!12},
  lab/.style={font=\scriptsize, inner sep=1.5pt},
  >={Stealth[length=1.6mm]}]
  % camada de procedência (coluna da esquerda)
  \node[prov] (doc) at (0,0)  {Documento};
  \node[prov] (chk) at (0,-2) {Chunk};
  % camada de conhecimento (grade 2 x 3)
  \node[ent] (pes) at (5.0,0)  {Pessoa};
  \node[ent] (org) at (10.3,0)  {Organizacao};
  \node[ent] (lin) at (5.0,-2) {Linguagem};
  \node[ent] (tec) at (10.3,-2) {Tecnologia};
  \node[ent] (hw)  at (10.3,-4) {Hardware};
  % relações de procedência e ligação entre camadas
  \draw[->] (doc) -- node[lab, left] (tc) {TEM\_CHUNK} (chk);
  \draw[->, dashed] (doc) -- node[lab, above] {SOBRE} (pes);
  \draw[->, dashed] (chk) -- node[lab, above] {MENCIONA} (lin);
  % relações da camada de conhecimento
  \draw[->] (pes) -- node[lab, above] {TRABALHOU\_EM} (org);
  \draw[->] (pes) -- node[lab, left] {CRIOU} (lin);
  \draw[->] (org) -- node[lab, right] (dv) {DESENVOLVEU} (tec);
  \draw[->] (tec) -- node[lab, above] {ESCRITO\_EM} (lin);
  \draw[->] (tec) -- node[lab, right] (ex) {EXECUTA\_EM} (hw);
  % caixas das camadas
  \begin{scope}[on background layer]
    \node[draw, dotted, inner sep=8pt, fit=(doc)(chk)(tc),
          label={[font=\scriptsize]above:Procedência}] {};
    \node[draw, dotted, inner sep=8pt, fit=(pes)(org)(hw)(dv)(ex),
          label={[font=\scriptsize]above:Conhecimento}] {};
  \end{scope}
\end{tikzpicture}%
}
\caption{Visão simplificada do modelo. Linhas tracejadas ligam as duas camadas; \rel{MENCIONA} e \rel{SOBRE} podem apontar para qualquer entidade. Os rótulos \rot{Premio} e \rot{Local} e parte das relações foram omitidos para simplificar.}
\label{fig:modelo}
\end{figure}

\subsection{Entidades e relações}

A camada de conhecimento tem sete tipos de entidade, descritos na Tabela~\ref{tab:entidades}. Além do rótulo específico, todo nó de entidade recebe o rótulo genérico \rot{Entidade}. Isso permite manter um único índice textual e um único índice vetorial para todas as entidades, de modo que a busca não precisa saber de antemão o tipo do que procura. Todas as entidades compartilham as propriedades \texttt{id}, \texttt{nome}, \texttt{aliases} (formas alternativas do nome), \texttt{descricao}, \texttt{wikidata\_id} e \texttt{embedding}.

\begin{table}[ht]
\centering
\caption{Tipos de entidade da camada de conhecimento.}
\label{tab:entidades}
\small
\begin{tabularx}{\textwidth}{lXl}
\toprule
\textbf{Rótulo} & \textbf{Representa} & \textbf{Propriedades específicas} \\
\midrule
\rot{Pessoa}      & Pesquisadores, engenheiros e fundadores & anos de nascimento e morte \\
\rot{Organizacao} & Empresas, universidades, laboratórios, agências e fundações & tipo, ano de fundação \\
\rot{Linguagem}   & Linguagens de programação & ano, paradigmas \\
\rot{Hardware}    & Computadores, processadores e periféricos & ano, tipo \\
\rot{Tecnologia}  & Sistemas, redes, protocolos, invenções e conceitos & ano, tipo \\
\rot{Premio}      & Prêmios e distinções & -- \\
\rot{Local}       & Países e cidades & tipo \\
\bottomrule
\end{tabularx}
\end{table}

Os rótulos \rot{Premio} e \rot{Local} não correspondem a páginas do corpus, mas foram incluídos porque aparecem com frequência nos textos e geram consultas multi-hop naturais, como ``quais ganhadores do Prêmio Turing trabalharam no Bell Labs?''.

Foram definidos 14 tipos de relação, cada um restrito a tipos específicos de origem e destino (Tabela~\ref{tab:relacoes}). Essas restrições funcionam como um esquema semântico: uma relação \rel{CRIOU}, por exemplo, só pode ligar uma pessoa a uma linguagem, um hardware ou uma tecnologia. Extrações que violam essas regras são rejeitadas antes de chegar ao banco, o que evita que erros do modelo de linguagem se tornem fatos no grafo.

\begin{table}[ht]
\centering
\caption{Tipos de relação e restrições de origem e destino.}
\label{tab:relacoes}
\small
\begin{tabularx}{\textwidth}{llX}
\toprule
\textbf{Relação} & \textbf{Origem} & \textbf{Destino} \\
\midrule
\rel{TRABALHOU\_EM} & Pessoa & Organizacao \\
\rel{ESTUDOU\_EM} & Pessoa & Organizacao \\
\rel{FUNDOU} & Pessoa & Organizacao \\
\rel{CRIOU} & Pessoa & Linguagem, Hardware, Tecnologia \\
\rel{DESENVOLVEU} & Organizacao & Linguagem, Hardware, Tecnologia \\
\rel{INFLUENCIADA\_POR} & Linguagem & Linguagem \\
\rel{ESCRITO\_EM} & Tecnologia & Linguagem \\
\rel{EXECUTA\_EM} & Tecnologia, Linguagem & Hardware \\
\rel{UTILIZA} & Hardware & Hardware \\
\rel{SUCESSOR\_DE} & qualquer & mesmo tipo da origem \\
\rel{RECEBEU} & Pessoa & Premio \\
\rel{COLABOROU\_COM} & Pessoa & Pessoa \\
\rel{NASCEU\_EM} & Pessoa & Local \\
\rel{SEDIADA\_EM} & Organizacao & Local \\
\bottomrule
\end{tabularx}
\end{table}

Toda relação extraída guarda três propriedades de procedência: \texttt{chunk\_ids}, a lista de trechos em que foi encontrada; \texttt{evidencia}, a frase do texto que a comprova; e \texttt{confianca}, um valor entre 0 e 1 atribuído na extração. Quando a mesma relação é encontrada em mais de um trecho, ela não é duplicada: o novo trecho é acrescentado à lista existente.

\subsection{Identificadores, restrições e índices}

Cada entidade recebe um identificador canônico formado pelo tipo e pelo nome normalizado, sem acentos e em minúsculas (por exemplo, \texttt{pessoa:alan\_turing} e \texttt{linguagem:cpp} para C++). Como todos os documentos usam a mesma regra, uma entidade citada em várias páginas converge para um único nó já na ingestão. Os casos que a regra não resolve, como ``Turing'' e ``Alan Turing'', são tratados na etapa de resolução de entidades (Seção~\ref{sec:ingestao}).

A integridade do grafo é garantida por 11 restrições de unicidade: sobre o \texttt{id} de cada um dos sete rótulos específicos, do rótulo \rot{Entidade}, de \rot{Documento} e de \rot{Chunk}, e sobre a URL de \rot{Documento}. Para as consultas, foram criados três tipos de índice, resumidos na Tabela~\ref{tab:indices}.

\begin{table}[ht]
\centering
\caption{Índices criados no Neo4j.}
\label{tab:indices}
\small
\begin{tabularx}{\textwidth}{lXX}
\toprule
\textbf{Tipo} & \textbf{Onde} & \textbf{Uso} \\
\midrule
Range (B-tree) & nome das entidades; categoria e título dos documentos; ano & Filtros e buscas exatas em consultas Cypher \\
Fulltext & nome, \texttt{aliases} e descrição das entidades; texto dos trechos & Busca por palavras-chave e por formas alternativas de nomes \\
Vetorial (cosseno) & \texttt{embedding} das entidades e dos trechos & Busca semântica a partir da pergunta do usuário \\
\bottomrule
\end{tabularx}
\end{table}

Os índices vetoriais têm 768 dimensões, correspondentes ao modelo de \emph{embedding} \texttt{multilingual-e5-base} \cite{wang2024me5}. O modelo foi escolhido por ser gratuito, executável localmente, treinado para vários idiomas, incluindo o português, e por aceitar entradas de até 512 \emph{tokens}, o que comporta os trechos do corpus sem truncamento.

\subsection{Verificação do esquema}

O esquema é criado por um script Cypher idempotente, que pode ser executado várias vezes sem duplicar estruturas. Um segundo script verifica o resultado: confere a existência das 11 restrições e dos índices (17 do tipo range, incluindo os criados automaticamente pelas restrições, 2 fulltext e 2 vetoriais, todos no estado \texttt{ONLINE}), insere nós de teste, executa uma consulta em cada índice textual e vetorial e remove os nós de teste ao final. Na versão 5.26 do Neo4j utilizada no projeto, todos os testes passaram, inclusive a busca textual por uma forma alternativa do nome armazenada na lista \texttt{aliases}.

% =====================================================================
% Extração de Entidades e Relações (Text-to-Graph)
% Responsável: Pessoa 2 (Anna Maria Albino)
%
% O que cobrir:
% - Como os textos foram divididos em chunks (tamanho, sobreposição) e por quê.
% - Qual LLM foi usado e com quais parâmetros (temperatura etc.).
% - O prompt de extração: como ele impõe os tipos do modelo (Tabelas 2 e 3) e as regras contra alucinação de docs/formato_extracao.md. Pode ir um trecho do prompt.
% - Formato de saída (JSON validado por src/contrato.py) e o que acontece com extrações inválidas.
% - Números: total de entidades e relações extraídas, % de JSONs aprovados pelo validador de primeira.
% - Avaliação manual numa amostra: quantas relações estavam corretas.
%
% Dicas: use \rot{Pessoa} para rótulos e \rel{CRIOU} para relações.
% Cite com \cite{chave} (chaves em referencias.bib). Ao terminar,
% apague o \pendente abaixo.
% =====================================================================

\section{Extração de Entidades e Relações (Text-to-Graph)}
\label{sec:extracao}

\pendente{extração de entidades e relações (text-to-graph)}

% =====================================================================
% Ingestão e Resolução de Entidades
% Responsável: Pessoa 3
%
% O que cobrir:
% - Pipeline: JSON validado -> nós e relações no Neo4j (uso de MERGE, constraints da Seção 4).
% - Geração dos embeddings com multilingual-e5-base (prefixo 'passage: ', normalização).
% - Resolução de entidades: como 'Turing' e 'Alan Turing' viram um nó só (aliases, fulltext, similaridade de embeddings, Wikidata).
% - Números: quantos nós/relações antes e depois da resolução, quantas fusões feitas.
% - Testes operacionais do pipeline (pytest).
%
% Dicas: use \rot{Pessoa} para rótulos e \rel{CRIOU} para relações.
% Cite com \cite{chave} (chaves em referencias.bib). Ao terminar,
% apague o \pendente abaixo.
% =====================================================================

\section{Ingestão e Resolução de Entidades}
\label{sec:ingestao}

\pendente{ingestão e resolução de entidades}

% =====================================================================
% Recuperação Híbrida e Raciocínio Multi-hop
% Responsável: Pessoa 4
%
% O que cobrir:
% - Busca vetorial (índices chunk_embedding e entidade_embedding, prefixo 'query: ').
% - Text-to-Cypher: como o LLM gera a consulta, como ela é validada antes de executar.
% - Como a busca vetorial e a consulta Cypher são combinadas (ordem, pesos, fusão de resultados).
% - Travessia multi-hop: exemplo de pergunta, a consulta Cypher gerada e o caminho percorrido.
% - Objeto de resposta enviado à interface (resposta, nós, arestas, Cypher, confiança).
%
% Dicas: use \rot{Pessoa} para rótulos e \rel{CRIOU} para relações.
% Cite com \cite{chave} (chaves em referencias.bib). Ao terminar,
% apague o \pendente abaixo.
% =====================================================================

\section{Recuperação Híbrida e Raciocínio Multi-hop}
\label{sec:recuperacao}

\pendente{recuperação híbrida e raciocínio multi-hop}

% =====================================================================
% Algoritmos de Grafo e Ancoragem das Respostas
% Responsável: Pessoa 5
%
% O que cobrir:
% - Algoritmos do Neo4j GDS usados (PageRank, Louvain/Leiden, Shortest Path) e para que cada um serve no contexto passado ao LLM.
% - Exemplo concreto: entidades mais centrais do grafo pelo PageRank, comunidades encontradas.
% - Grounding: como o LLM é obrigado a citar os caminhos do subgrafo e os trechos (evidências).
% - Como o nível de confiança da resposta é calculado.
%
% Dicas: use \rot{Pessoa} para rótulos e \rel{CRIOU} para relações.
% Cite com \cite{chave} (chaves em referencias.bib). Ao terminar,
% apague o \pendente abaixo.
% =====================================================================

\section{Algoritmos de Grafo e Ancoragem das Respostas}
\label{sec:grounding}

\pendente{algoritmos de grafo e ancoragem das respostas}

% =====================================================================
% Interface e Explicabilidade
% Responsável: Pessoa 6
%
% O que cobrir:
% - Arquitetura da interface em Streamlit e fluxo de uso.
% - Visualização do subgrafo (biblioteca usada) — incluir uma captura de tela como figura.
% - Como a procedência é exibida: nós consultados, Cypher gerado, trechos de evidência com link para a Wikipédia, confiança.
%
% Dicas: use \rot{Pessoa} para rótulos e \rel{CRIOU} para relações.
% Cite com \cite{chave} (chaves em referencias.bib). Ao terminar,
% apague o \pendente abaixo.
% =====================================================================

\section{Interface e Explicabilidade}
\label{sec:interface}

\pendente{interface e explicabilidade}

% =====================================================================
% Avaliação e Resultados
% Responsável: Pessoa 5
%
% O que cobrir:
% - Conjunto de perguntas de teste (quantas, quantos saltos cada uma). Ver exemplos em docs/modelo_lpg.md.
% - Comparação GraphRAG x RAG só vetorial: acertos, alucinações, respostas com fonte.
% - Qualidade das relações extraídas comparada ao Wikidata (usar os wikidata_id do dataset, Seção 3).
% - Tabela de resultados e discussão: onde o grafo ajudou, onde falhou.
%
% Dicas: use \rot{Pessoa} para rótulos e \rel{CRIOU} para relações.
% Cite com \cite{chave} (chaves em referencias.bib). Ao terminar,
% apague o \pendente abaixo.
% =====================================================================

\section{Avaliação e Resultados}
\label{sec:avaliacao}

\pendente{avaliação e resultados}

% =====================================================================
% Conclusão
% Responsável: Pessoa 1 (Maria Eduarda) — escrever depois da avaliação
%
% O que cobrir:
% - Retomar o objetivo e resumir o que foi construído.
% - Principais resultados, com números da Seção de Avaliação.
% - Limitações (ex.: páginas curtas, extração depende do LLM, só português).
% - Trabalhos futuros.
% =====================================================================

\section{Conclusão}

\pendente{conclusão}


\bibliographystyle{sbc}
\bibliography{referencias}

\end{document}
